\documentclass{article}
\usepackage{amsmath,amssymb}
\begin{document}
\title{A finite Laurent circuit flag for the biased butterfly decoder}
\maketitle

This is an exact finite algebra diagnostic, not a universality or density theorem.
The signed Haar decoder and its Farkas equations are defined in
\texttt{butterfly-biased-dual-audit.tex}.  Here $H=32$, the target is the first
saved target in \texttt{butterfly-adaptive-haar-h32-sample.json}, and
\[
 p_y(\varepsilon)=\sum_{j=0}^{4}y_j\varepsilon^{2^j},\qquad
 t_s=1-p_s.
\]
The original signed query matrix $A(\varepsilon)$ has $80$ rows and $63$
columns: one residue column per output and one bias column per input block.
Its row $e=(k,s,g)$ is
\[
 A_{e,y}=\frac{(1-2\theta_e)u_{k,g}(\pi^{-1}y)}{1-p_s-p_y},\qquad
 A_{e,(l,b)}=(1-2\theta_e)1_{(l,b)=(k,g)}.
\]

\paragraph{Exact calculation.}
At $\varepsilon=1/100$, rational row reduction gives rank $63$.  The selected
$63$-row square minor is consequently a nonzero rational function of
$\varepsilon$.  Gaussian elimination over rational Laurent series computes
all $17$ fundamental row dependencies.  The checker uses $1024$ initial
coefficients, a right-hand-side shift of $341$, and valuation-based pivoting.
Every division checks that the pivot's leading nonzero coefficient lies
within the available precision.  After normalizing and reducing the kernel
rows, at least $255$ coefficients remain known.  Every intermediate kernel
basis is also checked against every exact block-bias relation.

The file \texttt{butterfly-haar-circuit-flag-h32.json} records the target,
selected rows, pivot valuations, and the first nonzero coefficient of every
nonzero coordinate in the resulting kernel basis.  These are rational
series operations, not floating point fits.  The square minor is invertible
over the rational-function field, so the computed leading terms belong to
actual rational-function dependencies.  No identity is inferred merely
from a small numerical residual.

\paragraph{A limitation of a leading-support argument.}
The basis vector numbered $14$ from zero has leading coefficients $+1$ at
query rows $69$ and $77$, both at order zero.  All its other coordinates
have positive valuation.  Its first negative coordinates occur at rows
$75$ and $79$, with coefficient $-1$ at order $4$.
Thus the initial support is nonnegative, although the actual dependency is
not nonnegative for small positive $\varepsilon$.
An argument that excludes nonnegative circuits solely by the unrefined
initial row configuration cannot prove feasibility of this decoder.
It must retain later coefficients or use a different positive
transformation.  This observation does not produce a Farkas obstruction:
the later negative coordinates prevent that conclusion.  The target has a
separately saved positive-margin primal Laurent certificate.

\paragraph{Reproduction and scope.}
Run \texttt{heavy python3 artifacts/head-compression/}
\texttt{butterfly-haar-circuit-flag-check.py 5}.
The argument \texttt{4} checks a separate sixteen-token target.
The calculation is not Lean-checked and changes no model definition.
It establishes neither a bound on the density of feasible targets nor an
upper or lower bound for worst-case head complexity.
\end{document}
